\reading{CR-17}{Margin (Collateral) and Settlement}{GRIND}{4}

\begin{crkeybox}[Learning objectives for this reading]
\begin{enumerate}[label=\textbf{\alph*.},leftmargin=2em]
\item Describe the rationale for collateral management.
\item Describe the terms of a collateral agreement and features of a credit support annex (CSA) within the ISDA Master Agreement including threshold, initial margin, minimum transfer amount and rounding, haircuts, credit quality, and credit support amount.
\item Calculate the credit support amount (margin) under various scenarios.
\item Describe the role of a valuation agent.
\item Describe the mechanics of collateral and the types of collateral that are typically used.
\item Explain the process for the reconciliation of collateral disputes.
\item Explain the features of a collateralization agreement.
\item Differentiate between a two-way and one-way CSA agreement and describe how collateral parameters can be linked to credit quality.
\item Explain aspects of collateral including funding, rehypothecation, and segregation.
\item Explain how market risk, operational risk, and liquidity risk (including funding liquidity risk) can arise through collateralization.
\item Describe the various regulatory capital requirements.
\end{enumerate}
\end{crkeybox}

\begin{crtrapbox}[Single-layer reading]
CR-17 exists only in the Crash layer.
\end{crtrapbox}

% =====================================================================
\lo{CR-17 a}{Describe the rationale for collateral management.}

Collateral management is the process of handling collateral exchanges to manage
credit risk, often involving bilateral agreements where \emph{both} parties can be
required to post collateral.

\begin{itemize}
\item It reduces credit risk exposure, allowing for more trading.
\item It lowers capital requirements for firms.
\end{itemize}

\textbf{Evolution.} Collateral management transitioned from informal practices to
standardized processes with the introduction of ISDA documentation in 1994.
\term{Credit support annex (CSA)} agreements specify the terms for collateral
exchange to mitigate credit risk.

% =====================================================================
\lo{CR-17 b}{Describe the terms of a collateral agreement and features of a credit support annex (CSA) within the ISDA Master Agreement including threshold, initial margin, minimum transfer amount and rounding, haircuts, credit quality, and credit support amount.}

\begin{enumerate}
\item \term{Initial margin}, also called the \term{independent amount}. Defines an
amount of extra collateral that must be posted upfront, \emph{irrespective of the
exposure}.
\item \term{Threshold}. Defines the level of MtM above which collateral is posted.
When the exposure is above the threshold, the threshold amount is
\textbf{under-collateralized}. When the exposure is below the threshold it is
\textbf{not collateralized at all}.
\item \term{Minimum transfer amount (MTA)}. The minimum amount of collateral that
can be transferred. Exposure must exceed the \emph{sum} of threshold and MTA
before a collateral call can be made.
\item \term{Rounding}. A collateral call or return amount is always rounded to a
certain lot size, to avoid unnecessarily small amounts.
\item \term{Haircut}. A reduction in the value of the asset, to account for the
fact that its price may fall between the last collateral call and liquidation in
the event of default.
\item \term{Credit quality}. Collateral parameters can be linked to the
counterparty's credit quality; see objective h.
\item \term{Credit support amount}. The amount of margin that may be requested at
a given point in time; calculated at objective c.
\end{enumerate}

\begin{center}
\begin{tikzpicture}[font=\sffamily\footnotesize]
\draw[FRMGrey, line width=1pt, -{Stealth[length=6pt]}] (0,0) -- (13.6,0)
  node[right, font=\scriptsize\itshape, text=FRMGrey] {};
\node[below=16pt, font=\scriptsize\itshape, text=FRMGrey] at (6.8,0)
  {exposure (mark-to-market value), increasing};
% zones
\fill[FRMGreen!14] (0,0.12) rectangle (4.6,1.5);
\fill[FRMGold!16] (4.6,0.12) rectangle (8.4,1.5);
\fill[FRMRed!13] (8.4,0.12) rectangle (13.2,1.5);
\draw[FRMRule, line width=0.6pt] (0,0.12) rectangle (13.2,1.5);
\draw[FRMGrey, line width=0.8pt] (4.6,0.12) -- (4.6,1.5);
\draw[FRMGrey, line width=0.8pt] (8.4,0.12) -- (8.4,1.5);
\node[text width=4.2cm, align=center, font=\scriptsize] at (2.3,0.81)
  {\textbf{below the threshold}\\not collateralized at all};
\node[text width=3.5cm, align=center, font=\scriptsize] at (6.5,0.81)
  {\textbf{above threshold, below threshold $+$ MTA}\\no call can yet be made};
\node[text width=4.4cm, align=center, font=\scriptsize] at (10.8,0.81)
  {\textbf{above threshold $+$ MTA}\\call made, rounded to the lot size};
% markers
\draw[FRMNavy, line width=0.9pt] (4.6,-0.12) -- (4.6,0.12);
\node[below, font=\scriptsize\bfseries, text=FRMNavy] at (4.6,-0.14) {threshold $K$};
\draw[FRMNavy, line width=0.9pt] (8.4,-0.12) -- (8.4,0.12);
\node[below, font=\scriptsize\bfseries, text=FRMNavy] at (8.4,-0.14) {$K\;+\;$MTA};
\end{tikzpicture}
\end{center}
\figcap{The three zones a CSA creates. Two facts are easy to swap. First, the
threshold amount stays \emph{permanently} under-collateralized even once calls
begin: collateral is posted on the excess above $K$, not on the whole exposure.
Second, the trigger for a call is $K + \mathrm{MTA}$, not $K$ alone. Initial
margin sits outside this picture entirely, since it is posted irrespective of
exposure.}

% =====================================================================
\lo{CR-17 c}{Calculate the credit support amount (margin) under various scenarios.}

\begin{crdefbox}[The two quantities]
\term{Credit support amount.} ISDA CSA documentation defines this as the amount of
margin that may be requested at a given point in time. Because of parameters such
as thresholds, minimum transfer amounts and initial margins, \emph{it will not be
the same as the value of the portfolio}.

\smallskip
\term{Credit support balance.} The amount of margin held already.
\end{crdefbox}

\gapbadge

{\footnotesize\color{FRMGrey}This is a \emph{calculate} objective. The source notes
present two worked tables and then instruct the reader that there is ``no need to
get into above tables in very depth'', giving no formula. The formula below is
written from the GARP curriculum, and it reproduces both of the notes' tables
exactly.}

\begin{crgapbox}[What the objective asks for: the credit support amount formula]
Considering initially that a party is only \emph{receiving} margin, the relevant
amount considering only the threshold and initial margin is
\[ \max(\text{value} - K_c,\, 0) + \mathrm{IM}_c \]

More generally, including both parties, the equation defining the credit support
amount for \emph{variation} margin is
\[ \max(\text{value} - K_c,\, 0) \;-\; \max(-\text{value} - K_p,\, 0) \;-\; C \]

where value is the current value of the relevant transactions, $K_c$ and $K_p$ are
the thresholds for the counterparty and for the party making the calculation
respectively, $\mathrm{IM}_c$ is the initial margin for the counterparty, and $C$
is the credit support balance.

\smallskip
\textbf{Reading the sign.} A \emph{positive} result means margin can be called, or
requested to be returned. A \emph{negative} result indicates a requirement to post
or return margin. Either way this holds only as long as the amount is higher than
the minimum transfer amount, and rounding may then be applied.
\end{crgapbox}

\begin{crtrapbox}[Threshold and initial margin pull against each other]
The threshold and the initial margin act in \textbf{opposite directions}, and it
does not make sense to include both. Indeed, when initial margins are present the
threshold will usually be \emph{zero} --- which is exactly what the second worked
table below does.

\smallskip
Initial margin does not depend on any of the terms in the variation margin
equation and is computed \emph{separately}. Traditionally, variation margin is
captured in the CSA by the credit support amount, while initial margin is referred
to separately as the independent amount. Initial margin is \textbf{not netted}
against variation margin and may also be segregated. Where both parties post
initial margin, the initial margins themselves are not netted and are paid
separately.
\end{crtrapbox}

\begin{crexbox}[Margin calculation without initial margin]
\centering\footnotesize
\begin{tabular}{l r r}
\toprule
 & \thead{Time 1} & \thead{Time 2} \\
\midrule
Portfolio value & 250 & 225 \\
Threshold & 100 & 100 \\
Credit support balance & 0 & 150 \\
\midrule
\textbf{Credit support amount} & \textbf{150} & \textbf{--25} \\
\bottomrule
\end{tabular}
\tcblower
\textbf{Time 1.} $\max(250-100,\,0) - 0 = 150$. The credit support amount is the
value minus the threshold, which means the portfolio is only \emph{partially}
collateralized: 100 remains uncovered by construction.

\smallskip
\textbf{Time 2.} $\max(225-100,\,0) - 150 = 125 - 150 = -25$. Even though the
portfolio is uncollateralized by $225 - 150 = 75$, the credit support amount is
$-25$, which corresponds to \textbf{margin being returned}, up to the threshold.
\end{crexbox}

\begin{crexbox}[Margin calculation with initial margin]
\centering\footnotesize
\begin{tabular}{l r r}
\toprule
 & \thead{Time 1} & \thead{Time 2} \\
\midrule
Portfolio value & 250 & 275 \\
Threshold (set to zero) & 0 & 0 \\
Credit support balance & 0 & 250 \\
\midrule
\textbf{Credit support amount} & \textbf{250} & \textbf{25} \\
\midrule
Initial margin (independent amount) & 40 & 35 \\
\bottomrule
\end{tabular}
\tcblower
\textbf{Time 1.} $\max(250-0,\,0) - 0 = 250$. With the threshold set to zero the
whole exposure is collateralized.

\smallskip
\textbf{Time 2.} $\max(275-0,\,0) - 250 = 25$.

\smallskip
\textbf{Initial margin.} The 40 and 35 are \emph{not} part of either credit
support amount. They are computed separately as the independent amount and are not
netted against variation margin.
\end{crexbox}

\begin{crtrapbox}[Minimum transfer amounts make margin path-dependent]
MTAs can cause \term{path dependency} in margin calculations. Two portfolios can
have identical values at times $t_2$ and $t_3$ but different values at $t_1$, and
the credit support amount at $t_3$ will then differ, coming out either smaller or
larger.

\smallskip
The mechanism: the amount of credit support required at $t_3$ depends on the
credit support balance at $t_2$, which in turn depends on whether a call was
triggered at $t_1$. A call that was suppressed by the MTA at $t_1$ leaves a
different balance standing for every later date.
\end{crtrapbox}

% =====================================================================
\lo{CR-17 d}{Describe the role of a valuation agent.}

\begin{enumerate}[label=\alph*)]
\item \term{Responsibilities.} Calculating credit exposure, market values, credit
support amounts, and managing collateral delivery and return. More precisely: to
calculate the current value including the impact of netting; to calculate the
market value of margin previously posted and adjust it by the relevant haircuts;
to calculate the total uncollateralised exposure; and to calculate the credit
support amount.
\item \term{Who can act.} Larger entities may act as valuation agents, or
third-party agents may be used to prevent disputes.
\end{enumerate}

\begin{crkeybox}[For centrally cleared transactions]
Margin disputes are not relevant, because the CCP is the valuation agent.
\end{crkeybox}

% =====================================================================
\lo{CR-17 e}{Describe the mechanics of collateral and the types of collateral that are typically used.}

\begin{itemize}
\item \term{Mechanics.} Agreements specify currency, collateral types, delivery
timing, interest rates, and so on. Trades are marked to market and netted
regularly. \textbf{The party with negative MtM exposure posts collateral}, and
cash is the most common form.
\item \term{Types.} Cash, government securities, corporate bonds, letters of
credit, and equity.
\end{itemize}

% =====================================================================
\lo{CR-17 f}{Explain the process for the reconciliation of collateral disputes.}

Collateral disputes may arise due to the valuation and population of trades,
market data and market closing time, netting rules, and the valuation of
collateral previously posted. Managing a dispute involves three steps.

\begin{enumerate}
\item The disputing party notifies the counterparty of its intent to dispute the
exposure \textbf{by the end of the day following the collateral call}.
\item All \emph{undisputed} amounts are transferred, and the reason for the
dispute is identified.
\item For unresolved disputes, the parties request quotes from several market
makers, \textbf{usually four}, for the MtM value.
\end{enumerate}

Reconciling trades on a regular basis can minimize potential disputes.

% =====================================================================
\lo{CR-17 g}{Explain the features of a collateralization agreement.}

The features are the CSA parameters set out at objective b: initial margin,
threshold, minimum transfer amount, rounding, and haircuts. The zone diagram there
shows how the threshold and MTA interact to determine when a call is actually
made.

% =====================================================================
\lo{CR-17 h}{Differentiate between a two-way and one-way CSA agreement and describe how collateral parameters can be linked to credit quality.}

\begin{center}
\footnotesize
\begin{tabular}{L{2.4cm} L{5.9cm} L{6.2cm}}
\toprule
 & \thead{One-way CSA} & \thead{Two-way CSA} \\
\midrule
Who posts & Only \emph{one} party posts collateral, upon a specific event such as
  a downgrade &
  \emph{Both} parties post collateral \\
Who benefits & Benefits the receiving party; presents risk to the posting party &
  Benefits both parties \\
When used & When counterparties have \textbf{significantly different} credit risks &
  When counterparties have \textbf{similar} credit risk \\
\bottomrule
\end{tabular}
\end{center}

\subsection*{Linking CSA terms to credit quality}

\begin{itemize}
\item \term{Reduced workload} when the counterparty has strong credit.
\item \term{Increased protection} when the counterparty has weak credit, through
collateral required upon downgrade.
\end{itemize}

% =====================================================================
\lo{CR-17 i}{Explain aspects of collateral including funding, rehypothecation, and segregation.}

\begin{enumerate}[label=\alph*)]
\item \term{Substitution.} Posting an equivalent value of other eligible
collateral.
\item \term{Rehypothecation.} Transferring posted collateral \emph{on} to other
counterparties as collateral.
\item \term{Segregation.} Legally protecting posted collateral.
\end{enumerate}

\begin{crtrapbox}[Rehypothecation and segregation are opposites]
Rehypothecation lets the receiver reuse the collateral, which improves funding for
the receiver and increases risk for the poster. Segregation does the reverse: it
ring-fences the collateral for the poster and removes the funding benefit from the
receiver. A CSA that permits one is by construction weakening the other.
\end{crtrapbox}

% =====================================================================
\lo{CR-17 j}{Explain how market risk, operational risk, and liquidity risk (including funding liquidity risk) can arise through collateralization.}

Collateralization, while helpful in mitigating counterparty default risk,
introduces its own set of challenges.

\begin{center}
\footnotesize
\begin{tabular}{L{3.4cm} L{11.1cm}}
\toprule
\thead{Risk} & \thead{How it arises through collateralization} \\
\midrule
\term{Market risk} & Collateral value can fluctuate due to market movements.
  Delays in collateral posting or receiving create a \textbf{margin period of
  risk}. \\
\term{Operational risk} & Errors can occur in collateral handling: missed calls,
  failed deliveries, system failures. Strong controls are crucial --- clear
  agreements, robust IT systems, accurate valuations. \\
\term{Liquidity and liquidation risk} & Liquidating collateral incurs costs
  (bid-ask spread, selling costs). Large-scale liquidation can negatively impact
  price. \\
\term{Funding liquidity risk} & CSAs require funding readiness for collateral
  calls. Risk increases with illiquid markets and higher funding costs. \\
\term{Default risk} & Default of the collateral security itself reduces its value,
  and haircuts may not fully cover the loss. \\
\term{Foreign exchange risk} & Collateral in different currencies introduces FX
  risk. \\
\bottomrule
\end{tabular}
\end{center}
\figcap{Six risks, of which the objective names three. Note the shape of the
trade: collateralization removes counterparty default risk and replaces it with a
set of risks that are each individually smaller but none of which is zero.}

% =====================================================================
\lo{CR-17 k}{Describe the various regulatory capital requirements.}

Non-clearable OTC transactions require banks to hold \emph{higher} capital
requirements. In 2011 the G20 established bilateral collateral requirements for
non-clearable OTC derivatives.

These regulatory requirements cover both \term{variation} and \term{initial}
margins, with the goal of minimizing systemic risk and reducing potential
\term{regulatory arbitrage} between clearable and non-clearable OTC transactions.
